\subsection{extraction\_times\_list.txt}
%%
\subsubsection{Purpose}

This file contains a line for each extra result to display at chosen times that may be different from the storing times defined by \param{dtstor} parameter.

This option is activated by the presence of the file in the input folder (no specific parameter is required in the \emph{input.solv}).

The results are stored in new output files (one per chosen time) with the same format as results stored with the \param{dtstor} timestep. 

%%
\subsubsection{Format description}
 Each line contains the time in seconds from the beginning of the simulation, at which the results extraction is expected, following by a chosen name for this result.

Times do not need to be ordered in the file.

Lines starting with a hash are ignored as well as empty lines or space only lines.
If two lines give the same time, the second one is ignored. 

Note that the name following the time is mandatory.
%%
\paragraph*{Sample file}
\pfg{Répertoire d'exemple: mettre ce fichier dans distrib/file\_template/input/... }
With the following input file

\begin{verbatim}
#
# some extractions times for specific data extractions
#
3  extract\_time\_3

10 time\_10
  
3     extract\_time\_1
#20  extract\_time\_20
\end{verbatim}

the code will generate 2 output files with names:
\begin{verbatim}
    ResultsTime_p0000d00m03s000_extract_time_3.csv
    ResultsTime_p0000d00m10s000_time_10.csv
\end{verbatim}

Note that the real extraction time is the nearest time computed according to the current timestep and may be different from the chosen time.